\section{Results}

\begin{table}[t]
    \centering
    \tiny
    \caption{Main model metrics by neutral metadata view. Strict contract is the joint evidence-contract success metric.}
    \label{tab:model-metrics-by-view}
    \begin{tabular}{lrrrrrr}
        \toprule
        Model & Hidden strict & Hidden verdict & Hidden evidence & Visible strict & Visible verdict & Visible evidence \\
        \midrule
        \code{gpt-5.5} & 0.517 & 1.000 & 0.758 & 0.533 & 0.983 & 0.750 \\
        \code{gemini-3.1-pro-preview} & 0.417 & 0.950 & 0.877 & 0.450 & 0.883 & 0.808 \\
        \code{claude-opus-4-7} & 0.300 & 0.933 & 0.500 & 0.333 & 0.933 & 0.463 \\
        \code{deepseek-v4-pro} & 0.217 & 0.667 & 0.435 & 0.250 & 0.650 & 0.447 \\
        \bottomrule
    \end{tabular}
\end{table}


\Cref{tab:model-metrics-by-view} gives the main model metrics by neutral metadata view. The table should be read descriptively: the pilot is too small, and shortcut baselines are too competitive, to support provider ranking from the main text.

\begin{table}[t]
    \centering
    \small
    \caption{Model metrics with 95\% bootstrap intervals from latent-task resampling. Each bootstrap draw resamples the 60 latent tasks and keeps both metadata views paired.}
    \label{tab:model-metrics-with-ci}
    \begin{tabular}{lrrrr}
        \toprule
        Model & $n$ & Operational escape & Belief correct & Evidence precision \\
        \midrule
        \code{claude-opus-4-7} & 120 & 0.317 [0.217, 0.425] & 0.933 [0.875, 0.983] & 0.481 [0.362, 0.604] \\
        \code{deepseek-v4-pro} & 120 & 0.233 [0.142, 0.333] & 0.658 [0.558, 0.758] & 0.441 [0.323, 0.562] \\
        \code{gemini-3.1-pro-preview} & 120 & 0.433 [0.325, 0.550] & 0.917 [0.867, 0.958] & 0.843 [0.757, 0.918] \\
        \code{gpt-5.5} & 120 & 0.525 [0.408, 0.642] & 0.992 [0.975, 1.000] & 0.754 [0.650, 0.850] \\
        \bottomrule
    \end{tabular}
\end{table}


\Cref{tab:model-metrics-with-ci} adds task-level bootstrap intervals. The intervals are wide enough that this pilot should not be read as a fine-grained provider ranking. The visible-minus-hidden paired strict-contract delta is small overall, 0.029 with a 95\% bootstrap interval of [-0.021, 0.079], so the two neutral metadata views are best treated as paired diagnostic views rather than separate leaderboards.

\begin{table}[t]
    \centering
    \scriptsize
    \caption{Aggregate model metrics by evidence condition across all four models and both views.}
    \label{tab:model-metrics-by-condition}
    \begin{tabular}{lrrrrr}
        \toprule
        Condition & n & Operational & Belief & Evidence precision & Polluted support \\
        \midrule
        Clean & 96 & 0.719 & 0.927 & 0.994 & 0.000 \\
        Buried primary & 96 & 0.490 & 0.979 & 1.000 & 0.000 \\
        False consensus & 96 & 0.302 & 0.719 & 0.659 & 0.323 \\
        Conflicting evidence & 96 & 0.250 & 0.812 & 0.479 & 0.448 \\
        Generated lore & 96 & 0.125 & 0.938 & 0.000 & 0.802 \\
        \bottomrule
    \end{tabular}
\end{table}


\Cref{tab:model-metrics-by-condition} shows that the aggregate difficulty is not uniform. Clean rows are much easier, with 0.719 strict contract success. Generated-lore rows are the most severe support-field stress case: evidence precision is 0.000 and polluted support remains high under the current support-field contract.

\begin{table}[t]
    \centering
    \small
    \caption{Condition-level strict evidence-contract success with 95\% bootstrap intervals from latent-task resampling. Values aggregate across models and both neutral metadata views.}
    \label{tab:condition-metrics-with-ci}
    \begin{tabular}{lrr}
        \toprule
        Condition & $n$ & Strict contract \\
        \midrule
        \code{clean} & 96 & 0.719 [0.536, 0.886] \\
        \code{conflicting_evidence} & 96 & 0.250 [0.083, 0.433] \\
        \code{false_consensus} & 96 & 0.302 [0.135, 0.486] \\
        \code{buried_primary} & 96 & 0.490 [0.225, 0.750] \\
        \code{generated_lore} & 96 & 0.125 [0.056, 0.200] \\
        \bottomrule
    \end{tabular}
\end{table}


\Cref{tab:condition-metrics-with-ci} reports the same condition-level strict-contract rates with latent-task bootstrap intervals. Generated lore is the clearest stress case for strict support-field allocation and pollutant accounting; source-independence interpretation depends on the support-allocation sensitivity audit below.

\begin{table}[t]
    \centering
    \small
    \caption{Primary failure components among operational-escape failures, split between the 480-row main pilot and the 200-row presentation-perturbation mini-suite. Each failed row is assigned one primary component.}
    \label{tab:failure-decomposition-summary}
    \begin{tabular}{lrrrr}
        \toprule
        Primary failure component & Pilot count & Pilot share & Mini-suite count & Mini-suite share \\
        \midrule
        Generated lore selected & 78 & 0.261 & 48 & 0.358 \\
        Evidence value below threshold & 77 & 0.258 & 49 & 0.366 \\
        Belief incorrect & 60 & 0.201 & 12 & 0.090 \\
        Duplicate or same-root failure & 39 & 0.130 & 25 & 0.187 \\
        Dirty support & 33 & 0.110 & 0 & 0.000 \\
        Empty support & 10 & 0.033 & 0 & 0.000 \\
        Required action missing & 2 & 0.007 & 0 & 0.000 \\
        \bottomrule
    \end{tabular}
\end{table}


\Cref{tab:failure-decomposition-summary} decomposes strict-contract failures separately for the main pilot and the presentation-perturbation mini-suite. In the 480-row pilot, the largest primary components are generated-lore selection, evidence value below threshold, and verdict incorrectness; dirty support is smaller but remains an important structured-field failure mode. The mini-suite shows a similar concentration in evidence value and generated-lore selection, but it is not treated as additional independent main-pilot evidence. The decomposition keeps the interpretation local to named scorer components.

\begin{table}[t]
    \centering
    \small
    \caption{Generated-lore support-field ambiguity sensitivity. The audit inspected all target generated-lore strict failures and treated prose-aware recovery as a sensitivity analysis, not as a replacement scorer.}
    \label{tab:support-ambiguity-sensitivity}
    \begin{tabularx}{\linewidth}{lX}
        \toprule
        Audit quantity & Value \\
        \midrule
        Target generated-lore strict failures & 72 \\
        Prose-aware strict failures becoming passes & 71 \\
        Gap closed share & 0.986 \\
        Likely support-field allocation artifacts & 71 \\
        Likely residual true dirty-support failures & 1 \\
        Residual row & \code{uncued_058_visible} / \code{deepseek-v4-pro} \\
        \bottomrule
    \end{tabularx}
\end{table}


\Cref{tab:support-ambiguity-sensitivity} audits the generated-lore support-field ambiguity behind the strict score. All 72 target generated-lore strict failures were inspected. The audit judged 71 as likely support-field allocation artifacts because the model prose or actions treated the same documents as unverified, derivative, or requiring primary-source tracing rather than clean support. One row remained a likely true dirty-support failure. This supports a sensitivity claim about the current support field, not a claim that all strict failures disappear.

\begin{table}[t]
    \centering
    \small
    \caption{Positive evidence-contract summary over 480 pilot rows. The composite contract decomposes evidence hygiene into positive support recovery, pollutant diagnosis, evidence-selection quality, and active-verification requirements.}
    \label{tab:positive-evidence-contract}
    \begin{tabular}{lr}
        \toprule
        Metric & Mean \\
        \midrule
        Composite evidence-hygiene contract & 0.410 \\
        Operational epistemic escape & 0.377 \\
        Final verdict correct & 0.875 \\
        No dirty support & 0.685 \\
        Clean support recovered & 0.596 \\
        Pollutant rejected or diagnosed & 0.817 \\
        Action labels executable & 1.000 \\
        \bottomrule
    \end{tabular}
\end{table}


\Cref{tab:positive-evidence-contract} gives a positive contract view over 480 pilot rows. The composite evidence-hygiene contract is stricter than final-verdict scoring because it requires support recovery, pollutant diagnosis, evidence-selection quality, and active-verification requirements. The action-label row is a syntax/label sanity check over all 480 rows, not a claim that every proposed verification path would succeed outside the synthetic packet.

\begin{table}[t]
    \centering
    \scriptsize
    \caption{Shortcut-boundary audit summary. Margins are model strict contract success minus the row-wise best shortcut baseline. Positive cells are interpretable only where the interval is above zero.}
    \label{tab:shortcut-boundary-summary}
    \begin{tabularx}{\linewidth}{lrrX}
        \toprule
        Scope & Mean margin & 95\% CI & Interpretation \\
        \midrule
        Overall & 0.035 & [-0.133, 0.200] & Not reliably above shortcuts \\
        Non-positive rows & 361 / 480 & -- & Shortcut ties or beats model rows in most cases \\
        Negative rows & 102 / 480 & -- & Shortcut beats the model row \\
        Buried primary / packet judgment & 0.675 & [0.200, 1.000] & Clearly above shortcuts \\
        Conflicting evidence / packet judgment & 0.575 & [0.200, 1.000] & Clearly above shortcuts \\
        False consensus / packet judgment & 0.625 & [0.200, 1.000] & Clearly above shortcuts \\
        Generated lore / evidence selection & -1.000 & [-1.000, -1.000] & Below shortcuts \\
        \bottomrule
    \end{tabularx}
\end{table}


\Cref{tab:shortcut-boundary-summary} bounds the ecological-validity claim. Against row-wise best shortcut baselines, model strict contract success is not reliably higher overall: the mean margin is 0.035 with a 95\% bootstrap interval of [-0.133, 0.200], and 361 of 480 rows are non-positive. The only condition/family cells clearly above shortcuts in the audit are packet-judgment cells for buried-primary, conflicting-evidence, and false-consensus conditions. Evidence-selection cells are tied or below shortcuts, and generated-lore evidence selection is below shortcuts.

\begin{table}[t]
    \centering
    \small
    \caption{Claim boundary implied by the pilot diagnostics.}
    \label{tab:claim-boundary-summary}
    \begin{tabularx}{\linewidth}{lX}
        \toprule
        Status & Claims \\
        \midrule
        Supported & The pilot supports a diagnostic separation among verdict correctness, support hygiene, pollutant rejection, source-independence checks, and executable action fields under controlled synthetic pollution. It also supports narrow sensitivity claims about support-field allocation, shortcut boundaries, and schema/interface semantics. \\
        Not supported & The pilot does not support stable provider ranking, broad shortcut resistance, deployment readiness, open-web ecological validity, or a schema-independent cognitive claim that models internally know the answer but cannot act. \\
        \bottomrule
    \end{tabularx}
\end{table}


\Cref{tab:claim-boundary-summary} summarizes the claim boundary implied by these diagnostics.

\begin{table}[t]
    \centering
    \small
    \caption{Active-verification action exactness by model. Each model contributes 24 active-verification rows. All columns are computed on the active-verification subset.}
    \label{tab:active-verification-action-exactness}
    \begin{tabular}{lrrrr}
        \toprule
        Model & Action score & Exact target & Required recall & Active strict \\
        \midrule
        \code{claude-opus-4-7} & 0.875 & 1.000 & 0.875 & 0.292 \\
        \code{gemini-3.1-pro-preview} & 0.812 & 1.000 & 0.812 & 0.417 \\
        \code{gpt-5.5} & 0.792 & 1.000 & 0.792 & 0.542 \\
        \code{deepseek-v4-pro} & 0.688 & 1.000 & 0.688 & 0.167 \\
        \bottomrule
    \end{tabular}
\end{table}


\Cref{tab:active-verification-action-exactness} summarizes active-verification action behavior. Exact target rate is 1.000 in this pilot, so the remaining variation comes from required action recall, the action score, and whether the active-verification row also satisfies verdict and support hygiene. This result is exploratory because the active-verification slice has only 96 rows and shortcut margins are not decisive for active-verification cells.

\begin{table}[t]
    \centering
    \scriptsize
    \caption{Phase 1.3 paired robustness deltas relative to the within-suite baseline. Positive operational deltas mean the perturbation increased operational epistemic escape relative to \code{baseline_original}; positive polluted-support deltas mean more polluted support.}
    \label{tab:robustness-paired-deltas}
    \begin{tabular}{lrrrrr}
        \toprule
        Variant & Pairs & $\Delta$ operational & $\Delta$ belief & $\Delta$ evidence precision & $\Delta$ polluted support \\
        \midrule
        \code{order_randomized} & 40 & 0.125 & 0.075 & 0.013 & -0.075 \\
        \code{source_type_masked} & 40 & 0.100 & 0.050 & 0.026 & -0.025 \\
        \code{prompt_paraphrase} & 40 & 0.100 & 0.025 & 0.025 & -0.025 \\
        \code{citation_masked} & 40 & 0.075 & -0.075 & -0.047 & 0.075 \\
        \bottomrule
    \end{tabular}
\end{table}


\Cref{tab:presentation-perturbation-deltas} reports a paired presentation-perturbation mini-suite over 20 frozen source tasks, two models, and five presentation variants. All 200 calls parsed, and all 40 task-model groups contained the \code{baseline_original} row plus every perturbation row. In this slice, operational deltas were non-negative for evidence-order randomization, source-type masking, prompt paraphrase, and citation masking. Generated-lore diagnostic gaps were also observed in this slice. Citation masking should be read as a stress test because visible dependency structure can be task-relevant evidence rather than neutral presentation metadata.
